\section{AHSME 1985}
        \begin{enumerate}
        	\item (\textit{AHSME 1985}) If $2x+1=8$, then $4x+1=$


 $\mathrm{(A)\ } 15 \qquad \mathrm{(B) \ }16 \qquad \mathrm{(C) \  } 17 \qquad \mathrm{(D) \  } 18 \qquad \mathrm{(E) \  }19$


 

	\item (\textit{AHSME 1985}) In an arcade game, the "monster" is the shaded sector of a circle of radius $1$ cm, as shown in the figure. The missing piece (the mouth) has central angle $ 60\degree$. What is the perimeter of the monster in cm?


 
                    \begin{figure}[H]
                        \centering
                        \includegraphics[width=0.5\textwidth]{Problem/AHSME_I_1985_2_Problem_Figure_1.png}
                    \end{figure}
                    


 $\mathrm{(A)\ } \pi+2 \qquad \mathrm{(B) \ }2\pi \qquad \mathrm{(C) \  } \frac{5}{3}\pi \qquad \mathrm{(D) \  } \frac{5}{6}\pi+2 \qquad \mathrm{(E) \  }\frac{5}{3}\pi+2$


 

	\item (\textit{AHSME 1985}) In right $\triangle ABC$ with legs $5$ and $12$, arcs of circles are drawn, one with center $A$ and radius $12$, the other with center $B$ and radius $5$. They intersect the hypotenuse in $M$ and $N$. Then $MN$ has length 


 
                    \begin{figure}[H]
                        \centering
                        \includegraphics[width=0.5\textwidth]{Problem/AHSME_I_1985_3_Problem_Figure_1.png}
                    \end{figure}
                    


 $\mathrm{(A)\ } 2 \qquad \mathrm{(B) \ }\frac{13}{5} \qquad \mathrm{(C) \  } 3 \qquad \mathrm{(D) \  } 4 \qquad \mathrm{(E) \  }\frac{24}{5}$


 

	\item (\textit{AHSME 1985}) A large bag of coins contains pennies, dimes and quarters. There are twice as many dimes as pennies and three times as many quarters as dimes. An amount of money which could be in the bag is


 $\mathrm{(A)\ } \$306 \qquad \mathrm{(B) \ }  \$333 \qquad \mathrm{(C)\ } \$342 \qquad \mathrm{(D) \  }  \$348 \qquad \mathrm{(E) \  }  \$360$


 

	\item (\textit{AHSME 1985}) Which terms must be removed from the sum


 
 \[\frac{1}{2}+\frac{1}{4}+\frac{1}{6}+\frac{1}{8}+\frac{1}{10}+\frac{1}{12}\]


 if the sum of the remaining terms is to equal $1$?


 $\mathrm{(A)\ } \frac{1}{4}\text{ and }\frac{1}{8} \qquad \mathrm{(B) \ }\frac{1}{4}\text{ and }\frac{1}{12} \qquad \mathrm{(C) \  } \frac{1}{8}\text{ and }\frac{1}{12} \qquad \mathrm{(D) \  } \frac{1}{6}\text{ and }\frac{1}{10} \qquad \mathrm{(E) \  }\frac{1}{8}\text{ and }\frac{1}{10}$


 

	\item (\textit{AHSME 1985}) One student in a class of boys and girls is chosen to represent the class. Each student is equally likely to be chosen and the probability that a boy is chosen is $\frac{2}{3}$ of the probability that a girl is chosen. The ratio of the number of boys to the total number of boys and girls is


 $\mathrm{(A)\ } \frac{1}{3} \qquad \mathrm{(B) \ }\frac{2}{5} \qquad \mathrm{(C) \  } \frac{1}{2} \qquad \mathrm{(D) \  } \frac{3}{5} \qquad \mathrm{(E) \  }\frac{2}{3}$


 

	\item (\textit{AHSME 1985}) In some computer languages (such as APL), when there are no parentheses in an algebraic expression, the operations are grouped from right to left. Thus, $a \times b - c$ in such languages means the same as $a(b-c)$ in ordinary algebraic notation. If $a\div b-c+d$ is evaluated in such a language, the result in ordinary algebraic notation would be


 $\mathrm{(A)\ } \frac{a}{b}-c+d \qquad \mathrm{(B) \ }\frac{a}{b}-c-d \qquad \mathrm{(C) \  } \frac{d+c-b}{a} \qquad \mathrm{(D) \  } \frac{a}{b-c+d} \qquad \mathrm{(E) \  }\frac{a}{b-c-d}$


 

	\item (\textit{AHSME 1985}) Let $a,a',b,b'$ be real numbers with $a$ and $a'$ nonzero. The solution to $ax+b=0$ is less than the solution to $a'x+b'=0$ if and only if 


 $\mathrm{(A)\ } a'b < ab' \qquad \mathrm{(B) \ }ab' < a'b \qquad \mathrm{(C) \  } ab < a'b' \qquad \mathrm{(D) \  } \frac{b}{a} < \frac{b'}{a'} \qquad \mathrm{(E) \  }\frac{b'}{a'} < \frac{b}{a}$


 

	\item (\textit{AHSME 1985}) The odd positive integers $1, 3, 5, 7, \ldots$, are arranged into five columns continuing with the pattern shown on the right. Counting from the left, the column in which $1985$ appears is the


 
                    \begin{figure}[H]
                        \centering
                        \includegraphics[width=0.5\textwidth]{Problem/AHSME_I_1985_9_Problem_Figure_1.png}
                    \end{figure}
                    


 $\mathrm{(A)\ } \text{first} \qquad \mathrm{(B) \ }\text{second} \qquad \mathrm{(C) \  } \text{third} \qquad \mathrm{(D) \  } \text{fourth} \qquad \mathrm{(E) \  }\text{fifth}$


 

	\item (\textit{AHSME 1985}) An arbitrary circle can intersect the graph of $y = \sin x$ in


 $\mathrm{(A)\  } \text{at most }2\text{ points} \qquad \mathrm{(B)\  }\text{at most }4\text{ points} \qquad \mathrm{(C) \  } \text{at most }6\text{ points} \qquad \mathrm{(D) \ } \text{at most }8\text{ points}$
$\mathrm{(E) \  }\text{more than }16\text{ points}$


 

	\item (\textit{AHSME 1985}) How many distinguishable rearrangements of the letters in $CONTEST$ have both the vowels first? (For instance, $OETCNST$ is one such arrangement, but $OTETSNC$ is not.)


 $\mathrm{(A)\ } 60 \qquad \mathrm{(B) \ }120 \qquad \mathrm{(C) \  } 240 \qquad \mathrm{(D) \  } 720 \qquad \mathrm{(E) \  }2520$


 

	\item (\textit{AHSME 1985}) Let $p$, $q$ and $r$ be distinct prime numbers, where $1$ is not considered a prime. Which of the following is the smallest positive perfect cube having $n = pq^2r^4$ as a divisor?


 $\mathrm{(A)\ } p^8q^8r^8 \qquad \mathrm{(B) }\left(pq^2r^2\right)^3 \qquad \mathrm{(C) } \left(p^2q^2r^2\right)^3 \qquad \mathrm{(D) } \left(pqr^2\right)^3 \qquad \mathrm{(E) \  }4p^3q^3r^3$


 

	\item (\textit{AHSME 1985}) Pegs are put in a board $1$ unit apart both horizontally and vertically. A rubber band is stretched over $4$ pegs as shown in the figure, forming a quadrilateral. Its area in square units is


 
                    \begin{figure}[H]
                        \centering
                        \includegraphics[width=0.5\textwidth]{Problem/AHSME_I_1985_13_Problem_Figure_1.png}
                    \end{figure}
                    


 $\mathrm{(A)\ } 4 \qquad \mathrm{(B) \ }4.5 \qquad \mathrm{(C) \  } 5 \qquad \mathrm{(D) \  } 5.5 \qquad \mathrm{(E) \  }6$


 

	\item (\textit{AHSME 1985}) Exactly three of the interior angles of a convex polygon are obtuse. What is the maximum number of sides of such a polygon?


 $\mathrm{(A)\ } 4 \qquad \mathrm{(B) \ }5 \qquad \mathrm{(C) \  } 6 \qquad \mathrm{(D) \  } 7 \qquad \mathrm{(E) \  }8$


 

	\item (\textit{AHSME 1985}) If $a$ and $b$ are positive numbers such that $a^b = b^a$ and $b = 9a$, then the value of $a$ is


 $\mathrm{(A)\ } 9 \qquad \mathrm{(B) \ }\frac{1}{9} \qquad \mathrm{(C) \  } \sqrt[9]{9} \qquad \mathrm{(D) \  } \sqrt[3]{9} \qquad \mathrm{(E) \  }\sqrt[4]{3}$


 

	\item (\textit{AHSME 1985}) If $ A = 20 \degree$ and $ B = 25 \degree$, then the value of $\left(1+\tan A\right)\left(1+\tan B\right)$ is


 $\mathrm{(A)\ } \sqrt{3} \qquad \mathrm{(B) \ }2 \qquad \mathrm{(C) \  } 1+\sqrt{2} \qquad \mathrm{(D) \  } 2\left(\tan A+\tan B\right) \qquad \mathrm{(E) \  }\text{none of these}$


 

	\item (\textit{AHSME 1985}) Diagonal $DB$ of rectangle $ABCD$ is divided into three segments of length $1$ by parallel lines $L$ and $L'$ that pass through $A$ and $C$ and are perpendicular to $DB$. The area of $ABCD$, rounded to the one decimal place, is 


 
                    \begin{figure}[H]
                        \centering
                        \includegraphics[width=0.5\textwidth]{Problem/AHSME_I_1985_17_Problem_Figure_1.png}
                    \end{figure}
                    


 $\mathrm{(A)\ } 4.1 \qquad \mathrm{(B) \ }4.2 \qquad \mathrm{(C) \  } 4.3 \qquad \mathrm{(D) \  } 4.4 \qquad \mathrm{(E) \  }4.5$


 

	\item (\textit{AHSME 1985}) Six bags of marbles contain $18$, $19$, $21$, $23$, $25$ and $34$ marbles, respectively. One bag contains chipped marbles only. The other $5$ bags contain no chipped marbles. Jane takes three of the bags and George takes two of the others. Only the bag of chipped marbles remains. If Jane gets twice as many marbles as George, how many chipped marbles are there?


 $\mathrm{(A)\ } 18 \qquad \mathrm{(B) \ }19 \qquad \mathrm{(C) \  } 21 \qquad \mathrm{(D) \  } 23 \qquad \mathrm{(E) \  }25$


 

	\item (\textit{AHSME 1985}) Consider the graphs of $y = Ax^2$ and $y^2+3 = x^2+4y$, where $A$ is a positive constant and $x$ and $y$ are real variables. In how many points do the two graphs intersect?


 $\mathrm{(A) \ }\text{exactly }4 \qquad \mathrm{(B) \ }\text{exactly }2 \qquad$


 $\mathrm{(C) \  }\text{at least }1,\text{ but the number varies for different positive values of }A \qquad$


 $\mathrm{(D) \  }0\text{ for at least one positive value of }A \qquad \mathrm{(E) \ }\text{none of these}$


 

	\item (\textit{AHSME 1985}) A wooden cube with edge length $n$ units (where $n$ is an integer $>2$) is painted black all over. By slices parallel to its faces, the cube is cut into $n^3$ smaller cubes each of unit edge length. If the number of smaller cubes with just one face painted black is equal to the number of smaller cubes completely free of paint, what is $n$?


 $\mathrm{(A)\ } 5 \qquad \mathrm{(B) \ }6 \qquad \mathrm{(C) \  } 7 \qquad \mathrm{(D) \  } 8 \qquad \mathrm{(E) \  }\text{none of these}$


 

	\item (\textit{AHSME 1985}) How many integers $x$ satisfy the equation 
 \[\left(x^2-x-1\right)^{x+2} = 1?\]


 $\mathrm{(A)\ } 2 \qquad \mathrm{(B) \ }3 \qquad \mathrm{(C) \  } 4 \qquad \mathrm{(D) \  } 5 \qquad \mathrm{(E) \  }\text{none of these}$


 

	\item (\textit{AHSME 1985}) In a circle with center $O$, $AD$ is a diameter, $ABC$ is a chord, $BO = 5$ and $\angle ABO = \ \stackrel{\frown}{CD} \ = 60^{\circ}$. Then the length of $BC$ is


 
                    \begin{figure}[H]
                        \centering
                        \includegraphics[width=0.5\textwidth]{Problem/AHSME_I_1985_22_Problem_Figure_1.png}
                    \end{figure}
                    


 $\mathrm{(A)\ } 3 \qquad \mathrm{(B) \ }3+\sqrt{3} \qquad \mathrm{(C) \  } 5-\frac{\sqrt{3}}{2} \qquad \mathrm{(D) \  } 5 \qquad \mathrm{(E) \  }\text{none of the above}$


 

	\item (\textit{AHSME 1985}) If 
 \[x = \frac{-1+i\sqrt{3}}{2} \qquad\text{and}\qquad y = \frac{-1-i\sqrt{3}}{2},\] where $i^2 = -1$, then which of the following is not correct?


 $\mathrm{(A)\ } x^5+y^5 = -1 \qquad \mathrm{(B) \ }x^7+y^7 = -1 \qquad \mathrm{(C) \  } x^9+y^9 = -1 \qquad$


 $\mathrm{(D) \  } x^{11}+y^{11} = -1 \qquad \mathrm{(E) \  }x^{13}+y^{13} = -1$


 

	\item (\textit{AHSME 1985}) A non-zero digit is chosen in such a way that the probability of choosing digit $d$ is $\log_{10}{(d+1)}-\log_{10}{d}$. The probability that the digit $2$ is chosen is exactly $1/2$ the probability that the digit chosen is in the set


 $\mathrm{(A)\ } \{2,3\} \qquad \mathrm{(B) \ }\{3,4\} \qquad \mathrm{(C) \  } \{4,5,6,7,8\} \qquad \mathrm{(D) \  } \{5,6,7,8,9\} \qquad \mathrm{(E) \  }\{4,5,6,7,8,9\}$


 

	\item (\textit{AHSME 1985}) The volume of a certain rectangular solid is $8$ cm<sup>3</sup>, its total surface area is $32$ cm<sup>2</sup>, and its three dimensions are in geometric progression. The sums of the lengths in cm of all the edges of this solid is


 $\mathrm{(A)\ } 28 \qquad \mathrm{(B) \ }32 \qquad \mathrm{(C) \  } 36 \qquad \mathrm{(D) \  } 40 \qquad \mathrm{(E) \  }44$


 

	\item (\textit{AHSME 1985}) Find the least positive integer $n$ for which $\frac{n-13}{5n+6}$ is a non-zero reducible fraction.


 $\mathrm{(A)\ } 45 \qquad \mathrm{(B) \ }68 \qquad \mathrm{(C) \  } 155 \qquad \mathrm{(D) \  } 226 \qquad \mathrm{(E) \  }\text{none of these}$


 

	\item (\textit{AHSME 1985}) Consider a sequence $x_1,x_2,x_3,\dotsc$ defined by:

 \begin{align*}&x_1 = \sqrt[3]{3}, \\ &x_2 = \left(\sqrt[3]{3}\right)^{\sqrt[3]{3}},\end{align*}

and in general

 \[x_n = \left(x_{n-1}\right)^{\sqrt[3]{3}} \text{ for } n > 1.\]


 What is the smallest value of $n$ for which $x_n$ is an integer?


 $\mathrm{(A)\ } 2 \qquad \mathrm{(B) \ }3 \qquad \mathrm{(C) \  } 4 \qquad \mathrm{(D) \  } 9 \qquad \mathrm{(E) \  }27$


 

	\item (\textit{AHSME 1985}) In $\triangle ABC$, we have $\angle C = 3\angle A$, $a = 27$ and $c = 48$. What is $b$?


 
                    \begin{figure}[H]
                        \centering
                        \includegraphics[width=0.5\textwidth]{Problem/AHSME_I_1985_28_Problem_Figure_1.png}
                    \end{figure}
                    


 $\mathrm{(A)\ } 33 \qquad \mathrm{(B) \ }35 \qquad \mathrm{(C) \  } 37 \qquad \mathrm{(D) \  } 39 \qquad \mathrm{(E) \  }\text{not uniquely determined}$


 

	\item (\textit{AHSME 1985}) In their base $10$ representations, the integer $a$ consists of a sequence of $1985$ eights and the integer $b$ consists of a sequence of $1985$ fives. What is the sum of the digits of the base $10$ representation of the integer $9ab$?


 $\mathrm{(A)\ } 15880 \qquad \mathrm{(B) \ }17856 \qquad \mathrm{(C) \  } 17865 \qquad \mathrm{(D) \  } 17874 \qquad \mathrm{(E) \  }19851$


 

	\item (\textit{AHSME 1985}) Let $\left\lfloor x\right\rfloor$ be the greatest integer less than or equal to $x$. Then the number of real solutions to $4x^2-40\left\lfloor x\right\rfloor+51 = 0$ is


 $\mathrm{(A)\ } 0 \qquad \mathrm{(B) \ }1 \qquad \mathrm{(C) \  } 2 \qquad \mathrm{(D) \  } 3 \qquad \mathrm{(E) \  }4$


 

\end{enumerate}